% Complemento C4: Marcos de referencia, grillas y remuestreo
\capituloComplemento[4]{Marcos de referencia, grillas y remuestreo}\label{cap:marcos}

\section{La idea}

\subsection{Sistemas, marcos y grillas}\label{sec:sistemas-marcos}

Un sistema de coordenadas define reglas matemáticas para expresar posiciones. Un marco de
referencia materializa esas reglas en una época concreta mediante puntos, estaciones y
observaciones.

Cuando se cambia de marco pueden bastar traslaciones, rotaciones y escala. Si existen
deformaciones regionales, se utilizan grillas de corrección que asignan desplazamientos
variables según la ubicación. En datos ráster, el cambio de coordenadas obliga además a
remuestrear porque los centros de píxel no coinciden exactamente.

\begin{resultado}[Mapa mental de la materia]
Las primeras transformaciones enseñan a reconocer invariantes. Las afinidades amplían la
familia permitiendo deformaciones. La geometría diferencial lleva esas ideas a superficies
curvas, y la geodesia las aplica al elipsoide, las proyecciones y los marcos de referencia.
\end{resultado}

\begin{idea}[La composición como operación central]
Una operación muy importante es la composición: en general, para
construir un sistema de coordenadas es necesario que haya conmutatividad a la hora de componer
movimientos por coordenadas.
\end{idea}

\section{Construcciones, cuentas y ejemplos}

\subsection{Marcos de referencia y transformaciones geodésicas}

Un sistema de coordenadas indica cómo expresar posiciones; un marco de referencia materializa
ese sistema mediante convenciones, época y puntos o estaciones observadas. Cambiar de marco
puede exigir traslaciones, rotaciones, escala y correcciones espaciales.

\subsubsection{Transformaciones paramétricas}

Las transformaciones tridimensionales de semejanza suelen organizarse como
\[
  X_2=t+(1+\mu)RX_1,
\]
donde \(t\) es una traslación, \(\mu\) una pequeña variación de escala y \(R\) una rotación.
Las fórmulas de Molodensky permiten transformar coordenadas geodésicas considerando cambios
de datum y de elipsoide sin pasar necesariamente por una proyección plana.

\subsubsection{Transformaciones mediante grillas}

Una grilla NTv2 almacena correcciones \(\Delta\varphi,\Delta\lambda\) en nodos regulares. Para
un punto interior se interpolan las correcciones de los cuatro nodos vecinos, típicamente por
interpolación bilineal:
\[
\Delta(u,v)=
(1-u)(1-v)\Delta_{00}
+u(1-v)\Delta_{10}
+(1-u)v\Delta_{01}
+uv\Delta_{11}.
\]
Este enfoque representa deformaciones regionales que una transformación global de pocos
parámetros no puede modelar.

\subsection{Discretización, grillas y remuestreo}

Al pasar de coordenadas continuas a una grilla aparecen decisiones sobre resolución,
origen, orientación y regla de muestreo. Una transformación geométrica rara vez lleva los
centros de píxel exactamente a centros de otra grilla; por eso se requiere remuestreo.

\begin{center}
\begin{tabularx}{.9\textwidth}{>{\bfseries}l X}
\toprule
Método & Característica\\
\midrule
Vecino más próximo & conserva valores originales; apropiado para variables categóricas\\
Bilineal & interpola cuatro celdas; suaviza datos continuos\\
Bicúbico & usa un vecindario mayor; mayor suavidad y costo\\
\bottomrule
\end{tabularx}
\end{center}

El jacobiano de la transformación permite estudiar deformaciones locales y convertir áreas
y densidades entre sistemas de coordenadas.

\section{En profundidad}

\subsection{Conexión con la práctica de Agrimensura}

\begin{center}
\begin{tabularx}{\textwidth}{>{\bfseries}p{47mm} X}
\toprule
Concepto & Aplicación\\
\midrule
Coordenadas geodésicas \((\varphi,\lambda,h)\) & salida y procesamiento de observaciones GNSS\\
Base tangente y métrica & distancias y ángulos locales sobre el elipsoide\\
Radios \(M,N\) & variación geométrica con la latitud\\
Proyección conforme & coordenadas planas para cartografía y mensura\\
Transformaciones de marco & compatibilidad entre realizaciones geodésicas\\
Grillas de corrección & transformación regional de alta precisión\\
Remuestreo & ortomosaicos, modelos de elevación y datos ráster\\
Jacobiano & escala, deformación y propagación local\\
\bottomrule
\end{tabularx}
\end{center}

\begin{idea}[Estrategia de estudio acumulativa]
Cada nuevo tema debería incorporarse con la secuencia
\[
\text{definición}\rightarrow
\text{representación geométrica}\rightarrow
\text{fórmula}\rightarrow
\text{ejemplo numérico}\rightarrow
\text{implementación}.
\]
Para la implementación pueden usarse Python, NumPy, SciPy, SymPy, \texttt{pyproj},
GeographicLib y QGIS según el capítulo.
\end{idea}
